\documentclass[pdftex,english,oribibl]{llncs}

%% Spracheinstellungen laden
\usepackage[english]{babel}

%% Schriftart in der Ausgabe/Eingabe
\usepackage[T1]{fontenc}
\usepackage{textcomp}
\usepackage[latin1]{inputenc}

%% Zitate
\usepackage[numbers]{natbib}
\bibliographystyle{abbrvnat}
%\bibliographystyle{dinat}
%\bibliographystyle{plainnat}
%\bibliographystyle{splncs}
%% Similar to option "sectionbib" but \refname instead of \bibname
\makeatletter
\renewcommand\bibsection{\section*{\refname\@mkboth{\MakeUppercase{\refname}}{\MakeUppercase{\refname}}}}
\makeatother

%% Index
%\usepackage{makeidx}
%\makeindex

%% PDF Einstellungen
% muss nach natbib geladen werden!
\usepackage{nameref}
\usepackage{varioref}
\usepackage[pdfusetitle,pdftex,colorlinks]{hyperref}
\hypersetup{pdfborder={0 0 0}}
\hypersetup{bookmarksdepth=3}
\hypersetup{bookmarksopen=true}
\hypersetup{bookmarksopenlevel=1}
\hypersetup{bookmarksnumbered=true}
\usepackage{color}
\hypersetup{colorlinks=false}

%\usepackage[section]{tocbibind}

\makeatletter
\gdef\@keywords{}
\def\keywords#1{\gdef\@keywords{#1}}
\gdef\@subtitle{}
\def\subtitle#1{\gdef\@subtitle{#1}}

%% modified from llncs
\renewenvironment{abstract}{%
  \list{}{\advance\topsep by0.35cm\relax\small%
          \leftmargin=1cm%
          \labelwidth=\z@%
          \listparindent=\z@%
          \itemindent\listparindent%
          \rightmargin\leftmargin}%
          \item[\hskip\labelsep\bfseries\abstractname]}{%
  \if!\@keywords!\else{\item[~]\item[\hskip\labelsep\bfseries\keywordname]\@keywords}\fi%
  \endlist}

\AtBeginDocument{%
  \if!\@subtitle!\else\hypersetup{pdfsubject={\@subtitle}}\fi
  \if!\@keywords!\else\hypersetup{pdfkeywords={\@keywords}}\fi
}
\makeatother

% llncs hyperref fix
\makeatletter
\providecommand*{\toclevel@author}{0}
\providecommand*{\toclevel@title}{0}
\makeatother

%% Grafiken
\usepackage[pdftex]{graphicx}
\DeclareGraphicsExtensions{.pdf,.jpg,.png}
\usepackage{subfigure}

%% Mathe
\usepackage{amsmath}
\usepackage{amssymb}

%% Listings
\usepackage{listings}
\lstset{escapechar=\%, frame=tb, basicstyle=\footnotesize}

%% Sonstiges
\newcommand{\TODO}[1]{\par\textcolor{red}{#1}\marginpar{\textcolor{red}{TODO}}}
\newcommand{\TODOX}[1]{\textcolor{red}{#1}\marginpar{\textcolor{red}{TODO}}}
\pagestyle{plain}

% Keine "Schusterjungen"
\clubpenalty = 10000
% Keine "Hurenkinder"
\widowpenalty = 10000 \displaywidowpenalty = 10000

%%%%%%%%%%%%%%%%%%%%%%%%%%%%%%%%%%%%%%%%%%%%%%%%%%%%%%%%%%%%%%%%%%%%%%%%%%%%%%%
%%% BEGIN DOCUMENT
%%%%%%%%%%%%%%%%%%%%%%%%%%%%%%%%%%%%%%%%%%%%%%%%%%%%%%%%%%%%%%%%%%%%%%%%%%%%%%%
\title{Simulation-based resilience prediction of microservice architectures}
% \subtitle{My (optional) Subtitle}
\author{Samuel Beck, Johannes G�nth�r, Christoph Zorn}
\institute{University of Stuttgart\\Reliable Software Systems (RSS)\\70569 Stuttgart, Germany}


\begin{document}

\maketitle

\begin{abstract}
	Current software simulators are tailored towards one specific purpose of conservative software simulation. Given the success of these tools it would be useful to run these tools on microservice architectures. This paper will focus on the development of a simulator that can be used for microservice architectures.
\end{abstract}

\section{Introduction}\label{sec:introduction}

	Hier werden wir beschreiben was wir erreicht haben und wie wir unser Projekt angegangen sind. Es wird au�erdem erl�utert wie wir in diesem Paper vorgehen und was f�r Themengebiete genauer betrachtet werden.
	
	%This is my introduction. %\citet{Shaw2003WritingGoodSoftwareEngineeringResearchPapersMinitutorial} wrote a paper with %hints on how to write good software engineering research papers. By the way, this was an %example for using the \textit{natbib} command \texttt{\textbackslash{}citet\{\}}.
	
	%Section~\ref{sec:anotherSection} presents everything one must know. The conclusions follow %in Section~\ref{sec:conclusions}.

\section{Another Section}\label{sec:anotherSection}
	
	%\textit{AspectJ} can be used to weave cross-cutting concerns into Java programs %\citep{AspectJ2007}. By the way, this was an example for using the \textit{natbib} command %\texttt{\textbackslash{}citep\{\}}.
	
	%We will now demonstrate how to use subfigures (see Figure~\ref{fig:subfig}). %Figure~\ref{fig:circle} shows a circle. A star is displayed in Figure~\ref{fig:star}.
	
	%\begin{figure}
		% \centering
		%\subfigure[This is a circle.]{\label{fig:circle}
		 % \includegraphics[width=0.3\textwidth]{figures/template_circle.pdf}
		%}
		%\subfigure[This is a star.]{\label{fig:star}
		%  \includegraphics[width=0.3\textwidth]{figures/template_star.pdf}
		%}
		%\caption{A circle (a) and a star (b). Note that any caption ends with a full stop %character.}
		% \label{fig:subfig}
	%\end{figure}
  
\section{Existing Work and Tools/Research}

\subsection{Tools in Comparison}\label{subsec:toolsInComparison}
Spigo:\newline\newline
\begin{itemize}
	\item lightweight simulator for microservice architectures written in go
	\item can simulate over half a dozen of microservice structures (i.e. monolith, storage)
	\item can simulate a failure of a system during runtime (execution is fixed and static)
	\item 
\end{itemize}
\begin{tabular}{l|l}
	Advantages & disadvantages \\ \hline
	easy input & hard to get overview on code\\
	output metrics & metrics difficult to understand \\
	& workflow not very obvious \\
	chaos monkey & only 1 chaos monkey \\
\end{tabular}

Palladio/Simulizar:\newline\newline
\begin{itemize}
	\item analyizing self-adaptive systems (cloud-computing)
	\item 
\end{itemize}
\begin{tabular}{l|l}
	advantages & disadvanteges \\ \hline
\end{tabular}

\section{Simulator / Dokumentation}\label{sec:simulator}


\section{Conclusions}\label{sec:conclusions}

  These are my conclusions.

\bibliography{template}

\end{document}
